% !TeX program = xelatex

%%%%%%%%%%%%%%%%%%%%%%%%%%%%%%%%%%%%%%%%%%%%%%%%%%%%
% 第三章证明
%%%%%%%%%%%%%%%%%%%%%%%%%%%%%%%%%%%%%%%%%%%%%%%%%%%%


\chapter{面向完整恢复的DSC辅助编码计算结果传输方法研究}

\section{定理\ref{proposition1}的证明} \label{<appendix-proposition1>}
为了证明此定理，对于每个集合 $\Pi_i$，定义一个不带模运算的索引集如下
\begin{equation}
	\hat{\Pi}_i = \{ \alpha_{i-1} + 1, \cdots, \alpha_{i-1} + \lceil \eta_i n \rceil \}, \quad \forall i \in \mathcal{J}.
	\label{eq:hat_pi}
\end{equation}
不难验证集合 $\Pi_i$ 和集合 $\hat{\Pi}_i$ 之间的映射是一一对应的。
因此，要证明定理\ref{proposition1}，只需证明对于每个 $1 \leq j \leq m$，这 $r$ 个索引 $\{j, j+m, \cdots, j+(r-1)m\}$ 被包含在由(\ref{eq:hat_pi})定义的 $r$ 个不同的集合 $\hat{\Pi}_i$ 中；原因是，经过模$m$运算，这些索引中的每一个都将被映射为索引 $j$。
为此，观察到
\begin{align}
	\{j, j+m, \cdots, j+(r-1)m\} &\subset \{1, \cdots, rm\} 
	\subset \bigcup_{i=1}^{|\mathcal{J}|} \hat{\Pi}_i, \quad \forall j \in [m], \label{eq:subset_rel} 
\end{align}
其中最后一个包含关系归因于条件(\ref{eq:prop_cond1})和(\ref{eq:prop_cond2})。
由于集合 $\hat{\Pi}_i$ 是互斥的，且每个集合 $\hat{\Pi}_i$ 最多包含 $m$ 个元素（因为 $\eta_i \leq \frac{H(X)}{\log_2 q} = \frac{m}{n}$），(\ref{eq:subset_rel})表明这 $r$ 个索引确实被包含在至少$r$ 个不同的集合 $\hat{\Pi}_i$ 中。证毕。





\section{BP译码算法}\label{appendix-BP}
该算法将从长度为 $m$ 的伴随式序列 $\{S_{X_\upsilon}^{(j)}\}_{j=1}^m$ \upcite{1042242} 中解码出长度为 $n$ 的序列 $\{X_\upsilon^{(j)}\}_{j=1}^n$。
具体而言，在相应的 Tanner 图中，变量节点由 $\mathbf{X}_\upsilon^{1:n}$ 的元素组成，并在下文中为方便表示记为 $\{x_i\}_{i=1}^n$；
同样，校验节点由 $\mathbf{S}_{X_\upsilon}^{1:m}$ 的元素组成，并记为 $\{s_j\}_{j=1}^m$。
对于每个符号 $a \in GF(q)$，设 $M_{i \to j}(a)$ 为从变量节点 $x_i$ 发往校验节点 $s_j$ 的消息，$M_{j \to i}(a)$ 为从校验节点 $s_j$ 发往变量节点 $x_i$ 的消息。
基于伴随式的 BP 解码算法涉及以下四个步骤\upcite{1042242, 4697499, 4259766}：

初始化（Initialization）: 根据 $\{X_\upsilon^{(j)}\}_{j=1}^n$ 的分布，设定
\begin{equation}
	M_{i \to j}(a) = \text{Pr}\{x_i = a\}, \quad \forall a \in GF(q).
	\label{eq:bp_init}
\end{equation}
初始化之后，以下两个步骤将交替执行。

水平步骤（Horizontal step）: 对于每个校验节点 $s_j$，计算发往每个连接的变量节点 $x_i$ 的消息如下
\begin{align}
	M_{j \to i}(a) = \sum_{\{a_k\}_{k \in N(j)}} &\delta(a_i = a \text{ and } \sum_{k \in N(j)} H_{jk} a_k = s_j) \nonumber \\
	&\cdot \prod_{k \in N(j) \backslash i} M_{k \to j}(a_k), \quad \forall a \in GF(q),
	\label{eq:bp_horizontal}
\end{align}
其中当事件 $E$ 为真（假）时，$\delta(E)$ 等于一（零），$H_{jk}$ 是奇偶校验矩阵 $\mathbf{H}^{n \times m}$ 的元素，$N(j) \backslash i$ 表示除 $x_i$ 以外连接到 $s_j$ 的变量节点集合。

垂直步骤（Vertical Step）:  对于每个变量节点 $x_i$，计算发往每个连接的校验节点 $s_j$ 的消息如下
\begin{equation}
	M_{i \to j}(a) = \frac{1}{Z_{i \to j}} \text{Pr}\{x_i = a\} \prod_{k \in N(i) \backslash j} M_{k \to i}(a), \quad \forall a \in GF(q),
	\label{eq:bp_vertical}
\end{equation}
其中 $Z_{i \to j}$ 是一个归一化因子，确保概率之和为一，$N(i) \backslash j$ 表示除 $s_j$ 以外连接到 $x_i$ 的校验节点集合。

判决（Decision）: 做出解码判决如下
\begin{equation}
	\hat{X}_\upsilon^{(i)} = \arg \max_{a \in GF(q)} \text{Pr}\{x_i = a\} \prod_{j \in N(i)} M_{j \to i}(a), \quad \forall i \in [n].
	\label{eq:bp_decision}
\end{equation}
解码过程将重复进行，直到满足 (\ref{sx}) 或达到最大迭代次数。


\chapter{面向线性恢复的DSC辅助编码计算结果传输方法研究}

\section{嵌套线性码编码方案}\label{<nested linear coding>}
本节介绍文献\cite{lim2019towards}中提出的嵌套线性码编码方案在无损且无边信息条件下的特化形式，该方案是推论\ref{bcst}的编码实现基础。

考虑节点$k\in \mathcal{K}$分别以编码速率$\tilde{R}_{c,k}$和$\hat{R}_{c,k}$并使用两个索引$m_k\in[2^{n\tilde{R}_{c,k}}]$和$l_k\in[2^{n\hat{R}_{c,k}}]$。定义$\tilde{R}_{c,\max}:=\max\{\tilde{R}_{c,1},\cdots,\tilde{R}_{c,K}\}$，$\hat{R}_{c,\max}:=\max\{\hat{R}_{c,1},\cdots,\hat{R}_{c,K}\}$。定义$\text{m}_k(m_k),m_k\in[2^{n\tilde{R}_{c,k}}]$为$m_k$的$q$进制展开，每当$\tilde{R}_{c,k}<\tilde{R}_{c,\max}$时，将补零至长度为$\lceil n\tilde{R}_{c,\max}/\log(q)\rceil$。为了便于理解，本文后续假设对于所有速率，$\frac{n\tilde{R}_{c,k}}{\log(q)}$和$\frac{n\hat{R}_{c,k}}{\log(q)}$均为整数。由于矢量形式的$\mathbf{m}_k$和标量形式的$m_k$一一映射，因此后文中$m_k$将被用来表示其$q$进制展开向量$\mathbf{m}_k(m_k)$。定义由下面的步骤生成的$K$个码本的集合为一个$(2^{n\tilde{R}_{c,1}},\cdots,2^{n\tilde{R}_{c,K}}, 2^{n\hat{R}_{c,1}},\cdots,2^{n\hat{R}_{c,K}},q,n)$嵌套线性码。
设参数满足$\epsilon^{''}>\epsilon^{'}>\epsilon$。
\begin{itemize}
	\item[(1)] 码本生成：固定一个有限域$\mathbb{F}_q$和一个参数$\epsilon^{'}\in(0,1)$。随机生成一个$\kappa\times n$的矩阵$\mathbf{G}\in \mathbb{F}_q^{\kappa\times n}$，以及抖动向量$d_k^n \in \mathbb{F}_q^n,\forall k\in[K]$，其中$\kappa=\frac{n(\tilde{R}_{c,\max}+\hat{R}_{c,\max})}{\log (q)}$，且$\mathbf{G}$和$d_k^n$的每个元素都是根据有限域$\mathbb{F}_q$上的均匀分布独立随机生成的。

	对于每一个信源$k\in\mathcal{K}$，通过下式生成一个线性码$\mathcal{C}_k$，即
	\begin{equation}
		u_k^n(m_k,l_k):=u_k^n(\text{m}_k,\text{l}_k)=[\text{m}_k,\text{l}_k]\mathbf{G}\oplus d_k^n,
	\end{equation}
	其中，索引$m_k$的范围是$[2^{n\tilde{R}_{c,k}}]$，索引$l_k$的范围是$[2^{n\hat{R}_{c,k}}]$。根据这种构造方式，每个码字都是独立同分布且均匀分布的，即每个符号都服从有限域上的均匀分布，并且码字之间是相互独立的。

	\item [(2)] 编码：由于采用恒等映射$U_k=X_k$，对于每个信源$k\in \mathcal{K}$，编码器寻找一个索引对$(m_k,l_k)\in [2^{n\tilde{R}_{c,k}}]\times [2^{n\hat{R}_{c,k}}]$，使得生成的码字$u_k^n(m_k,l_k)$与信源序列$x_k^n$满足$\epsilon^{'}$-联合典型性，即
	\begin{equation}
		(u_k^n(m_k,l_k),x_k^n)\in \mathcal{T}_{\epsilon^{'}}^{(n)}.
	\end{equation}
	如果存在多个这样的索引对，则从中随机挑选一个。如果不存在这样的索引对，则从所有可能的索引对$[2^{n\tilde{R}_{c,k}}]\times [2^{n\hat{R}_{c,k}}]$中随机选择一个。节点$k$将发送索引$m_k$至移动设备。

	给定一个系数矩阵$\mathbf{A}\in \mathbb{F}_q^{J\times K}$，定义索引的线性组合为$\text{m}_\mathbf{A}(\text{m}_1,\cdots,\text{m}_K)=\mathbf{A}\begin{bmatrix}
			\text{m}_1\\\vdots \\\text{m}_K
	\end{bmatrix}$和$\text{l}_\mathbf{A}(\text{l}_1,\cdots,\text{l}_K)=\mathbf{A}\begin{bmatrix}
	\text{l}_1\\\vdots\\\text{l}_K
\end{bmatrix}$。本文后续将用$\text{m}_\mathbf{A}$和$\text{l}_\mathbf{A}$代替$\text{m}_\mathbf{A}(\text{m}_1,\cdots,$ $\text{m}_K)$和$\text{l}_\mathbf{A}(\text{l}_1,\\\cdots,\text{l}_K)$。进一步可以得到，$m_{\mathbf{A}_j}=\oplus_{k=1}^K A_{jk}m_k$以及$l_{\mathbf{A}_j}=\oplus_{k=1}^K A_{jk}l_k$，其中$\mathbf{A}_j$表示矩阵$\mathbf{A}$的第$j$行。

\item[(3)]译码：移动设备在接收到发送的索引$(\text{m}_1,\cdots,\text{m}_K)$后，译码器搜索一个唯一的索引元组$\hat{\text{l}}_{\mathbf{A}}$，使得所有的码字满足$\epsilon$-联合典型性（其中$\epsilon^{'}<\epsilon$），即
\begin{equation}
	\left(u_1^n(\text{m}_1, \hat{\text{l}}_1),\cdots,u_K^n(\text{m}_K, \hat{\text{l}}_K)\right)\in \mathcal{T}_{\epsilon}^{(n)},
\end{equation}
	其中，$\hat{\text{l}}_k\in [2^{n\tilde{R}_{c,k}}], \forall k\in \mathcal{K}$。如果译码器找到了唯一的索引元组，则将$\hat{w}_\mathbf{A}^n(m_\mathbf{A},\hat{l}_\mathbf{A})=[m_{\mathbf{A}},\hat{l}_\mathbf{A}]\mathbf{G}\oplus d_{\mathbf{A}}^n$作为目标线性函数的无损重构值，其中$m_{\mathbf{A}}$是接收到的索引的线性组合。否则，如果不存在这样的索引元组，或者存在多个这样的元组，则译码失败。
\end{itemize}

\section{定理\ref{<linear>}的证明}\label{<linear-proof>}
推论\ref{bcst}中的可达速率区域$\mathcal{R}_{lossless}$对辅助矩阵$\mathbf{E}$取并集，因此选取特定的$\mathbf{E}$可以得到$\mathcal{R}_{lossless}$的可达内界。以下分别取$\mathbf{E}=\mathbf{I}_K$和$\mathbf{E}=\mathbf{a}$两个特殊值，推导出两个可达边界。

（1）完整恢复边界$\mathcal{R}_{sw}$：
取辅助矩阵$\mathbf{E}=\mathbf{I}_K$，即$K\times K$单位矩阵，此时$\hat{\mathbf{g}}_{\mathbf{E}}^{(t)}=\mathbf{I}_K\tilde{\mathbf{g}}_{\mathcal{K}}^{(t)}=[\tilde{\mathbf{g}}_{\pi(1)}^{(t)},\cdots,\tilde{\mathbf{g}}_{\pi(K)}^{(t)}]$，即解码器的目标是完整恢复所有参与传输节点的编码梯度。此时$L_{\mathbf{E}}=K$，推论\ref{bcst}中的约束要求对所有满足$0\leq L_{\mathbf{C}}<K$的矩阵$\mathbf{C}$取交集。根据文献\cite{lim2019towards}，当$\mathbf{E}$为单位矩阵时，推论\ref{bcst}的可达速率区域退化为经典的SW区域。其直观解释如下：当$\mathbf{E}=\mathbf{I}_K$时，对于每一个行满秩矩阵$\mathbf{C}\in\mathbb{F}_q^{L_{\mathbf{C}}\times K}$（$L_{\mathbf{C}}$从$0$到$K-1$取遍所有值），条件熵$H(\hat{g}_{\mathbf{E}}^{(t)}|\hat{g}_{\mathbf{C}\mathbf{E}}^{(t)})=H(\tilde{g}_{\mathcal{K}}^{(t)}|\mathbf{C}\tilde{g}_{\mathcal{K}}^{(t)})$，且$|\mathcal{T}|=K-L_{\mathbf{C}}$。遍历所有$\mathbf{C}$和$\mathcal{T}$所产生的约束等价于对所有子集$\mathcal{T}\subseteq\mathcal{K}$的经典SW约束\upcite{lim2019towards}
\begin{align}
		\sum_{k\in\mathcal{T}}\tilde{R}_{c,k}>H\left(\{\tilde{g}_k^{(t)}\}_{k\in\mathcal{T}}|\{\tilde{g}_n^{(t)}\}_{n\notin\mathcal{T}}\right),\quad\forall \mathcal{T}\subseteq \mathcal{K}.
\end{align}
取$\mathcal{T}=\mathcal{K}$，和速率下界为
\begin{align}
	\sum_{k\in\mathcal{K}}\tilde{R}_{c,k}>H\left(\{\tilde{g}_k^{(t)}\}_{k\in\mathcal{K}}\right),
\end{align}
记该区域为$\mathcal{R}_{sw}$。

（2）线性恢复边界$\mathcal{R}_{linear}$：
取辅助矩阵$\mathbf{E}=\mathbf{a}$，即解码系数行向量，大小为$1\times K$，此时解码器的目标直接为$\hat{\mathbf{g}}_{\mathbf{a}}^{(t)}=\sum_{k\in\mathcal{K}}a_k\tilde{\mathbf{g}}_k^{(t)}=\mathbf{g}^{(t)}$。因此$L_{\mathbf{E}}=1$。由于推论\ref{bcst}中的约束条件(b)要求$0\leq L_{\mathbf{C}}<L_{\mathbf{E}}=1$，故$L_{\mathbf{C}}$只能取$0$，即$\mathbf{C}$为空矩阵。此时$|\mathcal{S}|=|\mathcal{T}|=L_{\mathbf{E}}-L_{\mathbf{C}}=1$。

由于$\mathcal{S}\subseteq\{1,\cdots,L_{\mathbf{E}}\}=\{1\}$，故$\mathcal{S}=\{1\}$。对于$\mathcal{T}$，由约束条件(d)要求
\begin{equation}
	\mathrm{rank}\left(\begin{bmatrix}
		\mathbf{E}(\mathcal{S})\\
		\mathbf{I}_K(\{1,\cdots,K\}\backslash\mathcal{T})
	\end{bmatrix}\right)=K.\label{full rank}
\end{equation}
其中$\mathbf{E}(\mathcal{S})=\mathbf{a}$为$1\times K$的非零行向量，$\mathbf{I}_K(\{1,\cdots,K\}\backslash\mathcal{T})$为$\mathbf{I}_K$中去掉$\mathcal{T}$对应行后的$(K-1)\times K$子矩阵。设$\mathcal{T}=\{k\}$，则$\mathbf{I}_K(\{1,\cdots,K\}\backslash\{k\})$张成除第$k$维外的所有维度。要使式(\ref{full rank})满秩，需要$\mathbf{a}$的第$k$个分量$a_k\neq 0$。由于梯度编码方案保证解码系数均为非零元素，故上述秩条件对所有$k\in\mathcal{K}$均成立。

由于$\mathbf{C}$为空，条件熵退化为无条件熵，对每个$\mathcal{T}=\{k\}$，代入推论\ref{bcst}的速率约束可得
\begin{align}
	\tilde{R}_{c,k}>H(\hat{g}_{\mathbf{a}}^{(t)})=H\left(\sum_{j\in\mathcal{K}}a_j\tilde{g}_j^{(t)}\right)=H(g^{(t)}),\quad\forall k\in \mathcal{K}.
\end{align}
其和速率下界为
\begin{align}
	\sum_{k\in\mathcal{K}}\tilde{R}_{c,k}\geq K\cdot H(g^{(t)}),
\end{align}
记该区域为$\mathcal{R}_{linear}$。

综上所述，$\mathcal{R}_{sw}$和$\mathcal{R}_{linear}$分别对应推论\ref{bcst}中$\mathbf{E}$取$\mathbf{I}_K$和$\mathbf{a}$时的可达内界。由于$\mathcal{R}_{lossless}$对$\mathbf{E}$取并集，两个特殊取值对应的区域之并$\mathcal{R}_{sub}=\mathcal{R}_{sw}\cup\mathcal{R}_{linear}$亦为$\mathcal{R}_{lossless}$的可达内界。证毕。

\section{推论\ref{condition}的证明}\label{<condition>}
考虑所有参与传输节点的信道条件相同的情况，即$h_{\pi(1)}=h_{\pi(2)}=\cdots=h_{\pi(K)}$。由于信道条件完全对称，优化问题$\mathcal{P}_{4.1}$的最优解具有对称结构，即各节点的最优信源编码速率相等，记为$\tilde{R}_{c,k}=\tilde{R}_c$，$\forall k\in\mathcal{K}$。此时传输时延$T_d$由各节点共同的信源编码速率$\tilde{R}_c$决定，因此比较两种策略下的最优$\tilde{R}_c$即可确定时延的优劣。

（1）对于$\mathcal{R}_{sw}$区域，将等速率$\tilde{R}_{c,k}=\tilde{R}_c$代入SW约束$\sum_{k\in\mathcal{T}}\tilde{R}_{c,k}>H(\{\tilde{g}_k^{(t)}\}_{k\in\mathcal{T}}|\{\tilde{g}_n^{(t)}\}_{n\notin\mathcal{T}})$，得到
\begin{align}
	\tilde{R}_c>\frac{H(\{\tilde{g}_k^{(t)}\}_{k\in\mathcal{T}}|\{\tilde{g}_n^{(t)}\}_{n\notin\mathcal{T}})}{|\mathcal{T}|},\quad\forall \mathcal{T}\subseteq \mathcal{K}.
\end{align}
因此，最优等速率为上式右端对所有$\mathcal{T}$取最大值。取$\mathcal{T}=\mathcal{K}$时，右端为$H(\{\tilde{g}_k^{(t)}\}_{k\in\mathcal{K}})/K$。对于梯度编码产生的信源，由于编码矩阵$\mathbf{B}$引入的线性相关性，可以验证$\mathcal{T}=\mathcal{K}$对应的约束最紧，即
\begin{align}
	\tilde{R}_c^{sw}=\frac{H(\{\tilde{g}_k^{(t)}\}_{k\in\mathcal{K}})}{K}.
\end{align}

（2）对于$\mathcal{R}_{linear}$区域，每个节点的速率约束为$\tilde{R}_{c,k}>H(g^{(t)})$，因此最优等速率为
\begin{align}
	\tilde{R}_c^{linear}=H(g^{(t)}).
\end{align}

当$\tilde{R}_c^{linear}<\tilde{R}_c^{sw}$时，即
\begin{align}
	H(g^{(t)})<\frac{H(\{\tilde{g}_k^{(t)}\}_{k\in\mathcal{K}})}{K},
\end{align}
等价于$K\cdot H(g^{(t)})<H(\{\tilde{g}_k^{(t)}\}_{k\in\mathcal{K}})$，线性恢复策略的每节点信源编码速率更低。由于信道条件相同，更低的信源编码速率直接对应更低的传输时延。证毕。

